\begin{homeworkProblem}[1-6]
  Seja $x : I \to \mathbb{R}^n$ uma solução de $x^{\prime}= f(x)$  com $0 \in I$. Mostre que as afirmações seguintes são equivalentes:
  \begin{enumerate}
    \item $x$ é constante.
    \item $x^{\prime}(t) = 0$ para todo $t \in I$.
    \item $x^{\prime}(0) = 0$.
    \item $f(x(0)) = 0$.
  \end{enumerate}
  Mostre ainda que se ocorrer uma dessas afirmações então $I=\mathbb{R}$.
  \begin{solution}
    É claro que se $x$ é constante, então $x^{\prime}(t)=0$ para todo $t\in I$.\\
    É claro que se $x^{\prime}(t)=0$ para todo $t\in I$, então $x^{\prime}(0)=0$.\\
    Note que se $x^{\prime}(0)=0$, usando que $x$ é uma solução de $x^{\prime}=f(x)$, temos que
    \begin{align*}
      f(x(0))=x^{\prime}(0)=0.
    \end{align*}
    Por último, note que $y(t)=x(0)$ para todo $t\in I$, temos que
    \begin{align*}
      y^{\prime}(t)=0=f(x(0))=f(y(t)).
    \end{align*}
    Pelo que temos que $y$ é solução de $x^{\prime}=f(x)$. Logo pela unicidade da solução, podemos garantir que $x$ é constante.\\
    Além disso, como $y(t)=x(0)$ para todo $t\in\mathbb{R}$, temos que $I=\mathbb{R}$.
  \end{solution}
\end{homeworkProblem}
